%! Intercepts, Zeros, and Extrema — Unit 2, Lesson 2.4 — Homework
\documentclass[10pt]{article}
\usepackage{atda-article}
\usepackage{atda-boxes}
\newcommand{\TallMath}[1]{$\displaystyle #1\rule[-1.4em]{0pt}{3.2em}$}
\setlength{\workrowsep}{8pt}

\begin{document}

\pageheader{Unit 2, Lesson 2.4}{Homework}
% NAMESTRIP: no \namedateperiod here. The student writes their name once, on the
% cover this component is stapled behind. See references/conventions.md.

\vspace{0.15in}

\boxguard[24]
\begin{notesbox}{Practice --- Everything Read Off a Graph}
\small
For every number you write, be ready to say which axis it came from.

\begin{enumerate}[leftmargin=1.6em, itemsep=6pt, topsep=4pt]

\item \textbf{Intercepts and zeros.} Fill the table from the three graphs below. Give intercepts as
points and zeros as numbers.

\begin{center}
\begin{tabular}{@{}c@{\hspace{0.4cm}}c@{\hspace{0.4cm}}c@{}}
\begin{tikzpicture}
\begin{axis}[
    width=5.0cm, height=3.9cm, title={\scriptsize (a) $y=x^2-1$},
    xmin=-3.4, xmax=3.4, ymin=-2.4, ymax=6.4,
    xtick={-3,-2,-1,0,1,2,3}, ytick={-2,0,2,4,6}, minor y tick num=1,
    grid=both, grid style={linegray!45}, minor grid style={linegray!30},
    axis lines=left, tick label style={font=\tiny}]
  \addplot[royal, very thick, domain=-2.6:2.6, samples=90]{x^2-1};
\end{axis}
\end{tikzpicture}
&
\begin{tikzpicture}
\begin{axis}[
    width=5.0cm, height=3.9cm, title={\scriptsize (b) $y=2^x-4$},
    xmin=-3.4, xmax=3.4, ymin=-5.4, ymax=4.4,
    xtick={-3,-2,-1,0,1,2,3}, ytick={-4,-2,0,2,4}, minor y tick num=1,
    grid=both, grid style={linegray!45}, minor grid style={linegray!30},
    axis lines=left, tick label style={font=\tiny}]
  \addplot[linegray, thick, dashed, domain=-3.3:3.3, samples=2]{-4};
  \addplot[royal, very thick, domain=-3.3:3, samples=90]{2^x-4};
\end{axis}
\end{tikzpicture}
&
\begin{tikzpicture}
\begin{axis}[
    width=5.0cm, height=3.9cm, title={\scriptsize (c) $y=2x+4$},
    xmin=-4.4, xmax=2.4, ymin=-4.4, ymax=8.4,
    xtick={-4,-3,-2,-1,0,1,2}, ytick={-4,-2,0,2,4,6,8}, minor y tick num=1,
    grid=both, grid style={linegray!45}, minor grid style={linegray!30},
    axis lines=left, tick label style={font=\tiny}]
  \addplot[royal, very thick, domain=-4.3:2.3, samples=2]{2*x+4};
\end{axis}
\end{tikzpicture}
\end{tabular}
\end{center}

\vspace{-2pt}
\renewcommand{\arraystretch}{1.4}
\begin{tabularx}{\linewidth}{@{}p{1.4cm} p{2.8cm} p{4.0cm} X@{}}
\rowcolor{mist}
\textbf{Graph} & \textbf{$y$-intercept} & \textbf{$x$-intercept(s)} & \textbf{Zero(s)} \\
\hline
(a) & \blank{2.0cm} & \blank{3.2cm} & \blank{2.8cm} \\
(b) & \blank{2.0cm} & \blank{3.2cm} & \blank{2.8cm} \\
(c) & \blank{2.0cm} & \blank{3.2cm} & \blank{2.8cm} \\
\end{tabularx}

\item \textbf{From an equation} (Lesson 2.3 spiral). For $y=(x-5)^2-4$: give the vertex, solve for
the two zeros, then give the $y$-intercept and the absolute minimum with its location.
\begin{work}
  y &= (x-5)^2-4 \\
  0 &= (x-5)^2-4 \\
  4 &= (x-5)^2 \\
  x-5 &= \pm 2 \\
  x &= 3 \text{ or } x = 7
\end{work}

Vertex: \blank{2.4cm} \quad $y$-intercept: \blank{2.4cm} \quad Absolute minimum: \blank{3.0cm}

\item \textbf{Input, output, or point?} Label each one.
\smallskip

\renewcommand{\arraystretch}{1.35}
\begin{tabularx}{\linewidth}{@{}X p{3.0cm}@{}}
\rowcolor{mist}
\textbf{The answer to\dots} & \textbf{is a\dots} \\
\hline
a zero of a function & \blank{2.6cm} \\
an $x$-intercept & \blank{2.6cm} \\
the \emph{value} of the absolute maximum & \blank{2.6cm} \\
the \emph{location} of the absolute maximum & \blank{2.6cm} \\
an interval where a function is decreasing & \blank{2.6cm} \\
\end{tabularx}

% Unguarded, item 4's stem and its (a)-(d) list sat at the foot of p1 with only one
% of its four write-line slots, the other three opening p2 -- a student's answer
% starting at the bottom of one page and finishing at the top of the next. \boxguard
% is inert inside a breakable tcolorbox, so \tcbbreak is the only lever. It cost
% nothing: still 4 pages, and p2 now carries items 4 and 5 plus the whole bike-share
% box through item 8. Mirrored in homework_key. See conventions.md ("boxguard").
\tcbbreak
\item \textbf{How many zeros?} For each function, say how many zeros it has and give the reason in a
few words. No graphing.
\begin{enumerate}[label=(\alph*), leftmargin=1.6em, itemsep=3pt, topsep=3pt]
  \item A parabola opening upward with vertex $(3,-4)$.
  \item A parabola opening upward with vertex $(3,4)$.
  \item An exponential whose graph creeps toward the line $y=0$ but never touches it.
  \item The horizontal line $y=5$.
\end{enumerate}

\writelines{3}

\item \textbf{The domain decides.} Consider $y=x^2$ on the restricted domain $2 \le x \le 5$. Give
the absolute maximum and the absolute minimum, each with its location, and say whether each sits at
a turning point or at an endpoint.

\writelines{2}

\end{enumerate}
\end{notesbox}

\vspace{0.10in}

\boxguard[30]
\begin{scenariobox}[In Context --- The Bike-Share Station]{royal}
\small
The graph shows $B(t)$, the number of bikes available at the station outside the school, $t$ hours
after $6$~a.m. The station holds twelve bikes when it is full.

\begin{center}
\begin{tikzpicture}
\begin{axis}[
    width=10.6cm, height=4.6cm,
    xlabel={$t$ (hours after 6 a.m.)}, ylabel={$B$ (bikes available)},
    xmin=0, xmax=12, ymin=0, ymax=13,
    xtick={0,1,2,3,4,5,6,7,8,9,10,11,12}, ytick={0,2,4,6,8,10,12},
    minor y tick num=1,
    grid=both, grid style={linegray!45}, minor grid style={linegray!30},
    axis lines=left,
    tick label style={font=\tiny}, label style={font=\scriptsize}]
  \addplot[royal, very thick] coordinates {(0,12) (3,0) (5,8) (8,8) (12,0)};
\end{axis}
\end{tikzpicture}
\end{center}

\begin{enumerate}[leftmargin=1.6em, itemsep=6pt, topsep=2pt, start=6]

\item Give the $y$-intercept as a point and both zeros of $B$ as numbers. Say what a zero means to
somebody walking up to this station.

\writelines{3}

\item Fill in the table from the graph, writing each interval in interval notation.
\smallskip

\renewcommand{\arraystretch}{1.4}
\begin{tabularx}{\linewidth}{@{}p{2.4cm} p{3.0cm} X@{}}
\rowcolor{mist}
\textbf{Behavior} & \textbf{On these inputs} & \textbf{What is happening at the station} \\
\hline
Decreasing & \blank{2.2cm} & \blank{6.2cm} \\
Increasing & \blank{2.2cm} & \blank{6.2cm} \\
Constant & \blank{2.2cm} & \blank{6.2cm} \\
Decreasing & \blank{2.2cm} & \blank{6.2cm} \\
\end{tabularx}

\item Give the absolute maximum and the absolute minimum of $B$, each as a value and a location.
One of the two happens at more than one location --- give every one of them.

\writelines{3}

\item The city wants to send a van to refill the station. Using the graph, say when you would send
it and justify the time with two characteristics you found above. Then settle this claim:
\emph{``The station ran out of bikes once during the day.''}

\writelines{4}

\end{enumerate}
\end{scenariobox}

\vspace{0.10in}

\boxguard[24]
\begin{scenariobox}[A First Look Ahead]{goldacc}
\small
\begin{enumerate}[leftmargin=1.6em, itemsep=6pt, topsep=2pt, start=10]

\item Coffee is poured into an insulated flask in a room kept at $20^\circ$C. The graph shows
$C(t)$, its temperature in degrees Celsius, $t$ hours after it was poured. The domain is every time
from the pour onward ($t \ge 0$); the graph shows the first six hours. Answer in plain English where
you do not have the vocabulary yet --- Lesson 2.5 will give you the words.

\begin{center}
\begin{tikzpicture}
\begin{axis}[
    width=10.0cm, height=4.2cm,
    xlabel={$t$ (hours after pouring)}, ylabel={$C$ (degrees C)},
    xmin=0, xmax=6, ymin=0, ymax=90,
    xtick={0,1,2,3,4,5,6}, ytick={0,20,40,60,80}, minor y tick num=1,
    grid=both, grid style={linegray!45}, minor grid style={linegray!30},
    axis lines=left,
    tick label style={font=\tiny}, label style={font=\scriptsize}]
  \addplot[linegray, thick, dashed, domain=0:6, samples=2]{20};
  \addplot[royal, very thick, domain=0:6, samples=100]{64*(0.5)^x+20};
\end{axis}
\end{tikzpicture}
\end{center}

\begin{enumerate}[label=(\alph*), leftmargin=1.6em, itemsep=5pt, topsep=4pt]
  \item Give the $y$-intercept as a point and say what it means about the coffee. Is $C$ increasing,
        decreasing, or constant, and on what interval?

\writelines{2}

  \item Does $C$ have any zeros? Answer yes or no and give the reason in terms of the coffee.

\writelines{1}

  \item Give the absolute maximum as a value and a location. Then decide whether $C$ has an absolute
        minimum, remembering the flask can sit there as long as you like.

\writelines{2}

  \item What number do the temperatures get closer and closer to, and does the graph ever reach it?
        (That flat dashed line has a name coming next lesson.)

\writelines{2}
\end{enumerate}

\end{enumerate}
\end{scenariobox}

\vspace{0.10in}

\boxguard[26]
\begin{extensionbox}
\small
\textbf{One domain change, three consequences.} Consider $y=x^2-4$.

\begin{enumerate}[label=(\alph*), leftmargin=1.6em, itemsep=5pt, topsep=4pt]
  \item On the domain of all real numbers, give the zeros, the $y$-intercept, and the absolute
        minimum with its location. Does it have an absolute maximum? Why or why not?

\writelines{3}

  \item Now restrict the domain to $1 \le x \le 4$ and nothing else changes. Complete the table.
\smallskip

\renewcommand{\arraystretch}{1.3}
\begin{tabular}{@{}|c|c|c|c|c|@{}}
\hline
$x$ & $1$ & $2$ & $3$ & $4$ \\ \hline
$y=x^2-4$ & \blank{1.2cm} & \blank{1.2cm} & \blank{1.2cm} & \blank{1.2cm} \\
\hline
\end{tabular}

\writelines{1}

  \item List the three things that changed when the domain was cut: how many zeros there are, where
        the absolute minimum is, and whether there is an absolute maximum. The rule $y=x^2-4$ was
        never touched --- so what does that tell you about the domain?

\writelines{3}
\end{enumerate}
\end{extensionbox}

\vspace{0.10in}

\begin{remindbox}
\small
\textbf{Next: Lesson 2.5 --- End Behavior and Asymptotes} (\texttt{AFDA.AF.2f--g}). Item 10 asked
what number the coffee's temperature creeps toward and whether it ever arrives, and the answer to
``does $y=2^x$ have a zero?'' has the same shape. Next lesson names that flat line an
\textbf{asymptote} and gives you language for what a graph does at the far left and far right of the
picture.
\end{remindbox}

\end{document}
